%% ============================================================
%%  Capítulo 2 – Arquitectura Final del Sistema
%% ============================================================
\chapter{Arquitectura Final del Sistema}
\label{chap:arquitectura}

\section{Descripción General}
\label{sec:arq_descripcion}

ALMA es un altavoz inteligente de mesa basado en arquitectura ARM64 que ejecuta Linux con
el stack de OpenVoiceOS (OVOS). El sistema integra procesamiento de voz en tiempo real,
reproducción de audio de alta fidelidad (24-bit/96\,kHz), interfaz gráfica táctil MIPI DSI e
iluminación ambiental RGB.

A nivel físico, el hardware de ALMA se distribuye en dos tarjetas de circuito impreso:
la \textbf{PCB-CTRL} (tarjeta principal de control) y la \textbf{PCB-PWR} (tarjeta de
gestión de energía). Esta partición, justificada en el Capítulo~\ref{chap:particion},
permite aislar los dominios de alta velocidad y señal analógica sensible de los circuitos de
conmutación de potencia.

% ── Plataforma de cómputo ────────────────────────────────────
\section{Plataforma de Cómputo Seleccionada}
\label{sec:arq_soc}

\begin{decisbox}[D-01 Cerrada: Selección de Plataforma de Cómputo]
  \textbf{Componente seleccionado:} Amlogic A311D2 (basado en la plataforma de
  referencia Khadas VIM4, 8\,GB).\\[4pt]
  \textbf{Justificación técnica:} El A311D2 es un SoC ARM64 big.LITTLE que incluye
  nativamente todos los periféricos requeridos por ALMA: interfaz MIPI DSI de 4 carriles,
  buses I2S y PDM para audio, I2C, SDIO y soporte para arranque seguro mediante TrustZone.
  El esquemático de referencia público de Khadas VIM4 permite validar el mapeo de pines y
  los circuitos de soporte, reduciendo el riesgo de errores de conexión.\\[4pt]
  \textbf{Impacto en PCB:} Encapsulado BGA de 1056 pines. Requiere un stack-up de mínimo
  6~capas y estrategia de escape HDI para el fan-out del BGA.
\end{decisbox}

\subsection{Configuración del SoC en Altium Designer}

Dada la complejidad del encapsulado BGA de 1056 pines, el componente fue segmentado en
\textbf{21 bloques esquemáticos (partes)} de aproximadamente 50 pines cada uno, organizados
secuencialmente desde los bancos de pines \texttt{A1} hasta \texttt{UA/UV}. Esta estrategia
sigue la referencia de Khadas y optimiza la legibilidad del esquemático.

\begin{infobox}[Mandato de Diseño: Pines tipo Pasivo]
  Todos los pines del A311D2 se configuraron como tipo \textbf{Pasivo} en la biblioteca
  de Altium Designer. Esta decisión es una estrategia deliberada de mitigación de riesgo:
  ante la ausencia de un modelo IBIS completo del fabricante, la configuración pasiva
  suprime los errores masivos de ERC (Electrical Rule Check) que se generarían por
  direccionalidad incorrecta de pines, trasladando la responsabilidad de validación
  a una revisión manual del netlist contra el esquemático de referencia de Khadas VIM4.
\end{infobox}

% ── Estrategia de audio ──────────────────────────────────────
\section{Estrategia de Audio y Micrófonos}
\label{sec:arq_audio}

\begin{decisbox}[D-02 Cerrada: Tipo de Micrófonos -- PDM Digital]
  \textbf{Selección:} Matriz de micrófonos MEMS digitales con interfaz PDM
  (Pulse Density Modulation).\\[4pt]
  \textbf{Justificación:} Los micrófonos PDM ofrecen mayor inmunidad al ruido
  electromagnético respecto a los analógicos con ADC externo, al trasladar la
  conversión A/D al interior del propio micrófono. Esto elimina la ruta de señal
  analógica en la PCB, reduce el riesgo de acoplamiento de ruido proveniente del
  amplificador Clase~D y del conmutador Buck, y simplifica el ruteo al tratarse de
  líneas digitales. El A311D2 dispone de interfaz PDM nativa, sin necesidad de
  hardware adicional.
\end{decisbox}

Para la salida de audio, se utiliza un \textbf{códec externo} conectado mediante bus
I2S, configurado para resolución de 24-bit a 96\,kHz. La configuración del códec se
realiza por bus I2C desde el SoC. El amplificador de potencia es de \textbf{Clase~D},
seleccionado por su eficiencia energética elevada ($>$85\,\%), reduciendo la potencia
disipada como calor dentro del enclosure \reqref{REQ-A-01}.



% ── Conectividad ─────────────────────────────────────────────
\section{Conectividad y Protocolos Soportados}
\label{sec:arq_conectividad}

\begin{decisbox}[D-04 Cerrada: Conectividad WiFi/BT -- Protocolo IP]
  \textbf{Selección:} Módulo WiFi dual-band 802.11ac (2.4/5\,GHz) con Bluetooth\,5.0 integrado,
  conectado al SoC mediante interfaz SDIO o USB.\\[4pt]
  \textbf{Protocolo de domótica:} ALMA opera exclusivamente como \textbf{cliente IP}
  en la red local. No actúa como hub de Zigbee, Thread ni Matter de forma directa;
  la interoperabilidad con ecosistemas domóticos se realiza vía red IP, interactuando
  con un hub doméstico existente. Toda comunicación saliente se cifra con TLS\,1.3
  \reqref{REQ-NF-02}.
\end{decisbox}

% ── Elemento de seguridad ────────────────────────────────────
\section{Elemento de Seguridad Criptográfica}
\label{sec:arq_seguridad}

\begin{decisbox}[D-05 Cerrada: TrustZone del A311D2]
  \textbf{Selección:} Se utiliza la implementación de ARM TrustZone integrada en
  el A311D2 como mecanismo principal de arranque seguro y gestión de claves.\\[4pt]
  \textbf{Justificación:} El TrustZone del A311D2 está documentado en el esquemático
  de referencia y no requiere componente adicional en PCB, reduciendo el BOM y el área
  de la tarjeta. Para el alcance del proyecto académico, este mecanismo cubre los
  13~controles mandatorios de ETSI\,EN\,303\,645 \reqref{REQ-E-02}. La adición de
  un Secure Element dedicado (p.\,ej. ATECC608B) se registra como decisión abierta
  para una eventual versión comercial.
\end{decisbox}

% ── Árbol de potencia ────────────────────────────────────────
\section{Árbol de Potencia del Sistema}
\label{sec:arq_potencia}

El sistema opera con una entrada nominal de \textbf{12\,V DC} provista por un adaptador
externo certificado bajo IEC\,62368-1 \reqref{REQ-E-01} medainte el protocolo de PD. La PCB-PWR convierte esta
tensión de entrada en los rieles regulados requeridos por cada subsistema. La
Tabla~\ref{tab:rieles_sistema} presenta el mapa completo de rieles; el detalle de
cálculos y componentes generadores se desarrolla en el Capítulo~\ref{chap:pcb_pwr}
y en el Anexo~\ref{chap:anexo_potencia}.

\begin{table}[H]
  \centering
  \caption{Mapa de rieles de alimentación del sistema ALMA.}
  \label{tab:rieles_sistema}
  \begin{tabularx}{\textwidth}{C{1.8cm} C{1.8cm} C{1.5cm} L{4cm} L{4cm}}
    \toprule
    \rowcolor{udea-green}
    \textcolor{white}{\textbf{Riel}} &
    \textcolor{white}{\textbf{Tensión}} &
    \textcolor{white}{\textbf{Tol.}} &
    \textcolor{white}{\textbf{Componente generador}} &
    \textcolor{white}{\textbf{Consumidores principales}} \\
    \midrule
    VIN    & 12\,V  & $\pm$5\,\%  & Conector USB C        & PCB-PWR (entrada) \\
    \rowcolor{row-alt}
    VDD\_NPU & 1.1\,V & $\pm$2\,\% & SY88003 (Silergy)       & NPU del A311D2 \\
    VDD\_IO  & 1.8\,V & $\pm$2\,\% & TL76701 (Texas Instr.)   & E/S SoC, buses lógicos \\
    \rowcolor{row-alt}
    VDD\_CPU & Variable & $\pm$2\,\% & MPU8756 (Mono. Power) & Clústeres CPU-A y CPU-B \\
    VCC5    & 5\,V   & $\pm$3\,\%  & NB679 (Mono. Power)      & Amplificador Clase~D, Boost pantalla \\
    \rowcolor{row-alt}
    VCC3V3  & 3.3\,V & $\pm$3\,\%  & NB680 (Mono. Power)     & WiFi/BT, eMMC, GPIO, sensores \\
    VLED    & 5\,V   & ---         & NB679 (compartido)       & LEDs RGB \\
    \bottomrule
  \end{tabularx}
\end{table}

\begin{riskbox}[Dispersión de Pines de Voltaje en el BGA]
  Los pines de alimentación del A311D2 (VDD\_CPU, VDD\_GPU, VDD\_NPU, etc.) están
  físicamente dispersos a lo largo del BGA de 1056 pines. Esta dispersión complica el
  diseño de los planos de potencia (Power Planes) en el layout y podría requerir capas
  adicionales dedicadas a distribución de potencia. Este riesgo se gestiona mediante
  un stack-up mínimo de 6~capas, con al menos dos capas internas dedicadas a GND y
  distribución de potencia (ver Capítulo~\ref{chap:fabricante}).
\end{riskbox}

% ── Organización de subsistemas ──────────────────────────────
\section{Organización de Subsistemas}
\label{sec:arq_subsistemas}

La Tabla~\ref{tab:subsistemas_pcb} resume la asignación de los seis subsistemas
identificados en el Entregable~2 a las dos tarjetas físicas del sistema.

\begin{table}[H]
  \centering
  \caption{Asignación de subsistemas a tarjetas PCB.}
  \label{tab:subsistemas_pcb}
  \begin{tabularx}{\textwidth}{C{1.5cm} L{3.5cm} C{2cm} L{5cm}}
    \toprule
    \rowcolor{udea-green}
    \textcolor{white}{\textbf{ID}} &
    \textcolor{white}{\textbf{Subsistema}} &
    \textcolor{white}{\textbf{PCB}} &
    \textcolor{white}{\textbf{Requisitos que satisface}} \\
    \midrule
    SPCA & Procesamiento Central y Almacenamiento & CTRL & \reqref{REQ-F-02}, \reqref{REQ-F-03}, \reqref{REQ-E-02} \\
    \rowcolor{row-alt}
    SAA  & Audio y Acústica           & CTRL & \reqref{REQ-F-01}, \reqref{REQ-F-02} \\
    SIVI & Interfaz Visual e Interacción & CTRL & \reqref{REQ-G-01}, \reqref{REQ-F-04} \\
    \rowcolor{row-alt}
    SGE  & Gestión de Energía         & PWR  & \reqref{REQ-E-01}, \reqref{REQ-A-01} \\
    SCR  & Conectividad y Red         & CTRL & \reqref{REQ-NF-02} \\
    \rowcolor{row-alt}
    SPSF & Privacidad y Seguridad Física & CTRL & \reqref{REQ-NF-01} \\
    \bottomrule
  \end{tabularx}
\end{table}
